\chapter{Les fractions comme nombres}\label{ch:g5:fractions}

Une fraction n'est pas seulement l'image d'une tarte~: c'est un
vrai \emph{nombre}, avec une place sur la droite graduée ---
parfois la même place qu'un nombre décimal. Ce chapitre relie les
deux mondes. (Les règles complètes du calcul fractionnaire arrivent
dans \cref{ch:g6:fractions} et au-delà.)

\section{Fractions sur la droite, encore}

\begin{example}[Placer et lire]\label{ex:g5:fractions:place}
Sur une droite découpée en cinquièmes, $\frac{7}{5}$ est à $7$ pas
de $0$~: après $1$ (qui est $\frac55$), à $1 + \frac25$. Chaque
fraction vit quelque part sur la droite~; un numérateur plus grand
(même dénominateur) signifie plus à droite.
\end{example}

\begin{proposition}[Même point, plusieurs noms]\label{prop:g5:fractions:equal}
Découper chaque part en deux (ou en trois\,\dots) ne déplace pas le
point~:
\[
\frac{1}{2} = \frac{2}{4} = \frac{3}{6} = \frac{5}{10},
\qquad
\frac{2}{3} = \frac{4}{6} = \frac{20}{30}.
\]
Multiplier (ou diviser) le numérateur et le dénominateur par le
même nombre donne un autre nom de la même fraction.
\end{proposition}
\admitted

\begin{omfigure}
\begin{tikzpicture}[scale=0.62]
  \foreach \i in {0,1} \draw[omExo] (\i*4,0) rectangle (\i*4+4,1.0);
  \fill[omDef!40] (0.06,0.06) rectangle (3.94,0.94);
  \draw[thick] (0,0) rectangle (8,1.0);
  \node[right, font=\small] at (8.3,0.5) {$\frac12$};
  \begin{scope}[yshift=-1.7cm]
  \foreach \i in {0,...,3} \draw[omExo] (\i*2,0) rectangle (\i*2+2,1.0);
  \fill[omDef!40] (0.06,0.06) rectangle (1.94,0.94);
  \fill[omDef!40] (2.06,0.06) rectangle (3.94,0.94);
  \draw[thick] (0,0) rectangle (8,1.0);
  \node[right, font=\small] at (8.3,0.5) {$\frac24$};
  \end{scope}
  \begin{scope}[yshift=-3.4cm]
  \foreach \i in {0,...,9} \draw[omExo] (\i*0.8,0) rectangle
    (\i*0.8+0.8,1.0);
  \foreach \i in {0,...,4} \fill[omDef!40] (\i*0.8+0.05,0.06)
    rectangle (\i*0.8+0.75,0.94);
  \draw[thick] (0,0) rectangle (8,1.0);
  \node[right, font=\small] at (8.3,0.5) {$\frac{5}{10}$};
  \end{scope}
\end{tikzpicture}

{\small Trois barres, trois noms, une même quantité~:
$\frac12 = \frac24 = \frac{5}{10}$.}
\end{omfigure}

\section{Fractions et décimaux}

\begin{proposition}[Fractions à nom décimal]\label{prop:g5:fractions:decimal}
Les fractions de dixièmes, centièmes, millièmes \emph{sont} des
nombres décimaux~:
\[
\frac{3}{10} = 0.3, \qquad \frac{27}{100} = 0.27 .
\]
Certaines autres fractions ont aussi un nom décimal, trouvé en
passant aux dixièmes ou aux centièmes~:
\[
\frac{1}{2} = \frac{5}{10} = 0.5, \qquad
\frac{1}{4} = \frac{25}{100} = 0.25, \qquad
\frac{3}{4} = 0.75, \qquad
\frac{1}{5} = \frac{2}{10} = 0.2 .
\]
Mais pas toutes~: $\frac13 = 0.333\dots$ ne s'arrête jamais --- la
fraction est son seul nom exact.
\end{proposition}
\admitted

\begin{omfigure}
\begin{tikzpicture}[scale=1.55]
  \draw[-Stealth] (-0.2,0) -- (7.6,0);
  \foreach \i in {0,...,28} \draw ({\i*0.25},-0.05) -- ({\i*0.25},0.05);
  \foreach \i in {0,4,...,28} \draw ({\i*0.25},-0.09) -- ({\i*0.25},0.09);
  \foreach \x in {0,...,7} \node[below, font=\small] at (\x,-0.1) {$\x$};
  \fill[omDef] (0.5,0) circle (1.6pt)
    node[above, font=\small] {$\frac12 = 0.5$};
  \fill[omThm] (2.75,0) circle (1.6pt)
    node[above, font=\small] {$\frac{11}{4} = 2.75$};
  \fill[omMeth] (6.25,0) circle (1.6pt)
    node[above, font=\small] {$\frac{25}{4} = 6.25$};
\end{tikzpicture}

{\small Une droite, deux langages~: chaque point peut porter un nom
fractionnaire et un nom décimal.}
\end{omfigure}

\begin{example}[Choisir le nom le plus commode]\label{ex:g5:fractions:choose}
Pour payer les trois quarts de $12$~: le nom fractionnaire est plus
commode, $\frac34$ de $12$ est $9$. Pour comparer $\frac34$ et
$0.8$~: le nom décimal est plus commode, $0.75 < 0.8$. Être
bilingue paie.
\end{example}

\section{Comparer des fractions}

\begin{method}[Trois outils de comparaison]\label{met:g5:fractions:compare}
\begin{enumerate}
  \item \emph{Même dénominateur}~: plus de parts gagne,
        $\frac57 > \frac37$~;
  \item \emph{même numérateur}~: les plus grandes parts gagnent,
        donc le \emph{plus petit} dénominateur gagne~:
        $\frac35 > \frac38$ (les cinquièmes sont plus grands que
        les huitièmes)~;
  \item \emph{repères}~: comparer chacune à $\frac12$ ou à $1$~:
        $\frac38 < \frac12 < \frac59$, donc $\frac38 < \frac59$.
\end{enumerate}
\end{method}

\begin{example}\label{ex:g5:fractions:compare}
Qui a mangé plus de pizza~: Ali avec $\frac58$ ou Bea avec
$\frac54$~? La fraction de Bea est plus grande que $1$ (cinq
quarts, c'est plus qu'une pizza entière), celle d'Ali est plus
petite que $1$~: Bea a mangé plus --- plus d'une pizza entière, en
fait.
\end{example}

\section{Exercices}

\begin{exercise}[$\star$]\label{exo:g5:fractions:1}
Dessiner une droite graduée découpée en cinquièmes de $0$ à $2$, et
placer~: $\frac25$~; \ $\frac55$~; \ $\frac85$~; \ $\frac{10}{5}$.
\end{exercise}

\begin{exercise}[$\star$]\label{exo:g5:fractions:2}
Compléter les noms égaux~:
$\frac12 = \frac{?}{8}$~; \quad
$\frac34 = \frac{?}{8}$~; \quad
$\frac{2}{5} = \frac{4}{?}$~; \quad
$\frac{6}{9} = \frac{2}{?}$.
\end{exercise}

\begin{exercise}[$\star$]\label{exo:g5:fractions:3}
Quelles fractions nomment des entiers~? Donner l'entier le cas
échéant~: $\frac82$~; \quad $\frac{9}{4}$~; \quad $\frac{15}{3}$~;
\quad $\frac{20}{5}$.
\end{exercise}

\begin{exercise}[$\star$]\label{exo:g5:fractions:4}
Écrire en décimaux~: $\frac{7}{10}$~; \quad $\frac{41}{100}$~; \quad
$\frac12$~; \quad $\frac34$~; \quad $\frac15$.
\end{exercise}

\begin{exercise}[$\star$]\label{exo:g5:fractions:5}
Écrire en fractions (dixièmes ou centièmes)~: $0.9$~; \quad
$0.13$~; \quad $2.5$~; \quad $0.75$ (puis donner son nom en
quarts).
\end{exercise}

\begin{exercise}[$\star$]\label{exo:g5:fractions:6}
Comparer, en nommant l'outil utilisé
(\cref{met:g5:fractions:compare})~:
\[
\frac47 \;?\; \frac67, \qquad
\frac35 \;?\; \frac37, \qquad
\frac25 \;?\; \frac{7}{12}, \qquad
\frac98 \;?\; \frac{11}{12} .
\]
\end{exercise}

\begin{exercise}[$\star$]\label{exo:g5:fractions:7}
Ranger avec les noms décimaux~: $\frac12$~; \ $0.4$~; \ $\frac34$~; \
$0.6$~; \ $\frac15$.
\end{exercise}

\begin{exercise}[$\star$]\label{exo:g5:fractions:8}
Un quart des élèves d'une classe de $28$ jouent d'un instrument, et
la moitié font du sport. Combien d'élèves pour chaque groupe~?
Lequel est le plus grand~?
\end{exercise}

\begin{exercise}[$\star$]\label{exo:g5:fractions:9}
Associer chaque fraction à une durée~: $\frac14$~h, $\frac12$~h,
$\frac34$~h, $\frac{1}{60}$~h --- à $30$~min, $45$~min, $1$~min,
$15$~min.
\end{exercise}

\begin{exercise}[$\star\star$]\label{exo:g5:fractions:10}
Tom a bu $\frac35$ de sa bouteille de $50$~cL~; Ana a bu $\frac12$
de sa bouteille de $60$~cL. Qui a bu le plus~? (Calculer les deux
quantités en cL --- comparer les fractions nues ne suffit pas ici,
et dire pourquoi.)
\end{exercise}

\begin{exercise}[$\star\star$]\label{exo:g5:fractions:11}
Donner une fraction strictement entre $\frac12$ et $1$~; puis une
strictement entre $\frac12$ et $\frac34$. (Les noms décimaux
peuvent aider.)
\end{exercise}
